% GENERATED FILE. Edit canonical lessons, then run npm run generate:books.
% canonical-source-hash: fnv1a64:3f12e0c7d24c1e35
% canonical-lessons: GU-C267-naraj, GU-C267-nirash, GU-C267-akalavu, GU-C267-cidavu, GU-C267-risavu

\chapter{Upset and disappointed}
\label{ch:gu-upset-and-disappointed}
\hlchaptermodality{267}

\begin{chapteropening}
\textbf{By the end of this chapter:} \emph{I can ask whether someone is upset and say that I am not disappointed.}

Complete the last lesson of chapter 267: I can ask whether someone is upset and say that I am not disappointed.
\end{chapteropening}

\section[\texorpdfstring{\gu{નારાજ} (nārāj)}{નારાજ (nārāj)}]{\gu{નારાજ} (nārāj) — displeased, upset}

\label{lesson:GU-C267-naraj}



\begin{quote}
Before the new one: say the Gujarati for to cry, then the Gujarati for to feel shy.
\end{quote}

\subsection*{You'll want to know: \gu{નારાજ}}
\textbf{\gu{નારાજ}} — \emph{nārāj} — \textquotedblleft{}displeased, upset\textquotedblright{}.

\begin{quote}
\textbf{\gu{શું} \gu{તમે} \gu{મારાથી} \gu{નારાજ} \gu{છો}?} — \emph{śũ tame mārāthī nārāj cho?} — \textquotedblleft{}Are you upset with me?\textquotedblright{}
\end{quote}

\begin{cousinweb}
\textbf{\gu{નારાજ}} (\emph{nārāj}) is a loan, usually traced to Persian \emph{nārāz}, \textquotedblleft{}displeased\textquotedblright{}. CDIAL has no entry for it.
\end{cousinweb}

\begin{grammarlens}[title={What you've built}]
One more word for everyday life at A2.
\end{grammarlens}

\subsection*{Your turn}
\begin{itemize}
\raggedright
  \item \emph{Say it:} \emph{nārāj}
  \item \emph{Say it:} \emph{nārāj}, once more
  \item \emph{Say it:} \emph{śũ tame mārāthī nārāj cho?} — Are you upset with me?
  \item \emph{Recall:} say the Gujarati for to cry, then the Gujarati for to feel shy, then say \emph{nārāj} again
\end{itemize}

\begin{tcolorbox}[breakable,colback=teal!4,colframe=teal!35!black,title={Before you move on}]
What does \gu{નારાજ} mean? (\textquotedblleft{}Displeased, upset\textquotedblright{}.) Say it once more.
\end{tcolorbox}
\section[\texorpdfstring{\gu{નિરાશ} (nirāś)}{નિરાશ (nirāś)}]{\gu{નિરાશ} (nirāś) — disappointed, hopeless}

\label{lesson:GU-C267-nirash}



\begin{quote}
Before the new one: say the Gujarati for to feel shy, then the Gujarati for displeased.
\end{quote}

\subsection*{You'll want to know: \gu{નિરાશ}}
\textbf{\gu{નિરાશ}} — \emph{nirāś} — \textquotedblleft{}disappointed, hopeless\textquotedblright{}.

\begin{quote}
\textbf{\gu{હું} \gu{પરિણામથી} \gu{નિરાશ} \gu{નથી}.} — \emph{hũ pariṇāmthī nirāś nathī} — \textquotedblleft{}I am not disappointed with the result.\textquotedblright{}
\end{quote}

\begin{cousinweb}
\textbf{\gu{નિરાશ}} (\emph{nirāś}) is Sanskrit \emph{nirāśa}, \textquotedblleft{}hopeless\textquotedblright{}, built from \emph{nir-}, \textquotedblleft{}without\textquotedblright{}, and \emph{āśā}, \textquotedblleft{}hope\textquotedblright{}, and taken over as a tatsama; CDIAL lists Gujarati \emph{nirās} (CDIAL 7278).
\end{cousinweb}

\begin{grammarlens}[title={What you've built}]
One more word for everyday life at A2.
\end{grammarlens}

\subsection*{Your turn}
\begin{itemize}
\raggedright
  \item \emph{Say it:} \emph{nirāś}
  \item \emph{Say it:} \emph{nirāś}, once more
  \item \emph{Say it:} \emph{hũ pariṇāmthī nirāś nathī} — I am not disappointed with the result.
  \item \emph{Recall:} say the Gujarati for to feel shy, then the Gujarati for displeased, then say \emph{nirāś} again
\end{itemize}

\begin{tcolorbox}[breakable,colback=teal!4,colframe=teal!35!black,title={Before you move on}]
What does \gu{નિરાશ} mean? (\textquotedblleft{}Disappointed, hopeless\textquotedblright{}.) Say it once more.
\end{tcolorbox}
\section[\texorpdfstring{\gu{અકળાવું} (akḷāvũ)}{અકળાવું (akḷāvũ)}]{\gu{અકળાવું} (akḷāvũ) — to feel restless, to be fed up}

\label{lesson:GU-C267-akalavu}



\begin{quote}
Before the new one: say the Gujarati for displeased, then the Gujarati for disappointed.
\end{quote}

\subsection*{You'll want to know: \gu{અકળાવું}}
\textbf{\gu{અકળાવું}} — \emph{akḷāvũ} — \textquotedblleft{}to feel restless, to be fed up\textquotedblright{}.

\begin{quote}
\textbf{\gu{શું} \gu{તમે} \gu{ગરમીમાં} \gu{અકળાઓ} \gu{છો}?} — \emph{śũ tame garmīmā̃ akḷāo cho?} — \textquotedblleft{}Do you feel restless in the heat?\textquotedblright{}
\end{quote}

\begin{cousinweb}
\textbf{\gu{અકળાવું}} (\emph{akḷāvũ}) is a tadbhava. CDIAL lists it, glossed \textquotedblleft{}to be perplexed\textquotedblright{}, under Sanskrit \emph{ākula}, \textquotedblleft{}confused; crowded\textquotedblright{} (CDIAL 1012).
\end{cousinweb}

\begin{grammarlens}[title={What you've built}]
One more word for everyday life at A2.
\end{grammarlens}

\subsection*{Your turn}
\begin{itemize}
\raggedright
  \item \emph{Say it:} \emph{akḷāvũ}
  \item \emph{Say it:} \emph{akḷāvũ}, once more
  \item \emph{Say it:} \emph{śũ tame garmīmā̃ akḷāo cho?} — Do you feel restless in the heat?
  \item \emph{Recall:} say the Gujarati for displeased, then the Gujarati for disappointed, then say \emph{akḷāvũ} again
\end{itemize}

\begin{tcolorbox}[breakable,colback=teal!4,colframe=teal!35!black,title={Before you move on}]
What does \gu{અકળાવું} mean? (\textquotedblleft{}To feel restless, to be fed up\textquotedblright{}.) Say it once more.
\end{tcolorbox}
\section[\texorpdfstring{\gu{ચિડાવું} (ciḍāvũ)}{ચિડાવું (ciḍāvũ)}]{\gu{ચિડાવું} (ciḍāvũ) — to get irritated, to get cross}

\label{lesson:GU-C267-cidavu}



\begin{quote}
Before the new one: say the Gujarati for disappointed, then the Gujarati for to feel restless.
\end{quote}

\subsection*{You'll want to know: \gu{ચિડાવું}}
\textbf{\gu{ચિડાવું}} — \emph{ciḍāvũ} — \textquotedblleft{}to get irritated, to get cross\textquotedblright{}.

\begin{quote}
\textbf{\gu{મારા} \gu{પિતા} \gu{ક્યારેય} \gu{નથી} \gu{ચિડાતા}.} — \emph{mārā pitā kyārey nathī ciḍātā} — \textquotedblleft{}My father never gets irritated.\textquotedblright{}
\end{quote}

\begin{cousinweb}
\textbf{\gu{ચિડાવું}} (\emph{ciḍāvũ}) is a tadbhava. It comes from \emph{*ciḍ}, \textquotedblleft{}be angry\textquotedblright{}, a reconstructed form; CDIAL glosses the Gujarati verb \textquotedblleft{}to be provoked\textquotedblright{} (CDIAL 4794).
\end{cousinweb}

\begin{grammarlens}[title={What you've built}]
One more word for everyday life at A2.
\end{grammarlens}

\subsection*{Your turn}
\begin{itemize}
\raggedright
  \item \emph{Say it:} \emph{ciḍāvũ}
  \item \emph{Say it:} \emph{ciḍāvũ}, once more
  \item \emph{Say it:} \emph{mārā pitā kyārey nathī ciḍātā} — My father never gets irritated.
  \item \emph{Recall:} say the Gujarati for disappointed, then the Gujarati for to feel restless, then say \emph{ciḍāvũ} again
\end{itemize}

\begin{tcolorbox}[breakable,colback=teal!4,colframe=teal!35!black,title={Before you move on}]
What does \gu{ચિડાવું} mean? (\textquotedblleft{}To get irritated, to get cross\textquotedblright{}.) Say it once more.
\end{tcolorbox}
\section[\texorpdfstring{\gu{રિસાવું} (risāvũ)}{રિસાવું (risāvũ)}]{\gu{રિસાવું} (risāvũ) — to sulk, to take offence}

\label{lesson:GU-C267-risavu}



\begin{quote}
Before the new one: say the Gujarati for to feel restless, then the Gujarati for to get irritated.
\end{quote}

\subsection*{You'll want to know: \gu{રિસાવું}}
\textbf{\gu{રિસાવું}} — \emph{risāvũ} — \textquotedblleft{}to sulk, to take offence\textquotedblright{}.

\begin{quote}
\textbf{\gu{શું} \gu{તે} \gu{ફરી} \gu{રિસાય} \gu{છે}?} — \emph{śũ te pharī risāy che?} — \textquotedblleft{}Is she sulking again?\textquotedblright{}
\end{quote}

\begin{cousinweb}
\textbf{\gu{રિસાવું}} (\emph{risāvũ}) is a tadbhava. It comes from Sanskrit \emph{ríṣyati}, \textquotedblleft{}is hurt, suffers wrong\textquotedblright{}; CDIAL glosses the Gujarati verb \textquotedblleft{}to be offended\textquotedblright{} (CDIAL 10749).
\end{cousinweb}

\begin{grammarlens}[title={What you've built}]
One more word for everyday life at A2.
\end{grammarlens}

\subsection*{Your turn}
\begin{itemize}
\raggedright
  \item \emph{Say it:} \emph{risāvũ}
  \item \emph{Say it:} \emph{risāvũ}, once more
  \item \emph{Say it:} \emph{śũ te pharī risāy che?} — Is she sulking again?
  \item \emph{Recall:} say the Gujarati for to feel restless, then the Gujarati for to get irritated, then say \emph{risāvũ} again
\end{itemize}

\begin{tcolorbox}[breakable,colback=teal!4,colframe=teal!35!black,title={Before you move on}]
What does \gu{રિસાવું} mean? (\textquotedblleft{}To sulk, to take offence\textquotedblright{}.) Say it once more.
\end{tcolorbox}
